\chapter{УДИРТГАЛ}
\begin{sloppy}

\section{Судалгааны үндэслэл ба асуудлын тодорхойлолт}

Монгол Улсын аялал жуулчлалын салбар нь байгаль, түүх, биет болон биет бус соёлын өвөөр баялаг хэдий ч эдгээр нөөцийг дижитал хэлбэрээр тайлбарлах, хэрэглэгчийн байршил ба нөхцөл байдалтай уялдуулан туршлага болгон хувиргах ажил хангалтгүй хэвээр байна. Уламжлалт танилцуулгын хэлбэр болох танилцуулгын самбар, хэвлэмэл гарын авлага, аман тайлбар нь мэдээлэл дамжуулах суурь үүргийг гүйцэтгэдэг ч хэрэглэгчийн оролцоо, сонирхлын ялгаа, хэлний олон янз байдал, хөдөлгөөнт нөхцөлийг хангалттай дэмждэггүй. Дижитал аялал жуулчлалын өнөөгийн чиг хандлагад байршилд суурилсан үйлчилгээ, ухаалаг гар утасны мэдрэгч, өргөтгөсөн бодит байдал, тоглоомжуулалтын механикийг хослуулсан интерактив туршлага илүү үр дүнтэй болохыг судалгаанууд харуулж байна \cite{Dias2021MobileTourism,Yovcheva2014AugmentedTourism,Xu2017GamificationTourism}.

Аялал жуулчлалын бүтээгдэхүүнийг зөвхөн мэдээлэл түгээх хэрэгсэл гэж үзэх бус, харин тухайн газарт очих, орчноо ажиглах, тайлбар нээх, соёлын утгыг урамшуулалтайгаар эзэмших үйл явц болгон загварчлах шаардлага үүсч байна. Ялангуяа залуу үеийн жуулчин, дотоодын аялагчид нь мобиль аппликейшн, тоглоомчилсон сорилт, персоналчилсан дижитал туршлагад илүү идэвхтэй ханддаг \cite{Jo2025ARGamification,Thinnukool2025LannaPassport}. Иймээс Монголын аялал жуулчлалын хэрэглээнд зориулсан AR суурьтай тоглоомын систем хөгжүүлэх нь зөвхөн техникийн сонирхолтой санаа бус, салбарын дижитал шилжилттэй шууд холбоотой судалгааны асуудал мөн.

Энэ ажлын төв асуудал нь дараах байдлаар тодорхойлогдоно: GPS-ийн нарийвчлал хязгаарлагдмал, мобайл төхөөрөмжийн мэдрэгчийн чанар харилцан адилгүй, орчны гэрэлтүүлэг ба хөдөлгөөн тогтворгүй нөхцөлд аялал жуулчлалын байршлыг тоглоомчилсон AR туршлагатай хэрхэн найдвартай, ойлгомжтой, академик үндэслэлтэй уялдуулах вэ? Энэхүү асуудлыг зөв шийдэхийн тулд онолын суурь, системийн зохиомж, хэрэгжилт, туршилтыг нэг мөр логик урсгалтай авч үзэх шаардлагатай.

\section{Судалгааны зорилго ба зорилтууд}

Энэхүү дипломын ажлын зорилго нь Монголын аялал жуулчлалын байршлуудыг илүү сонирхолтой, оролцоонд суурилсан байдлаар танилцуулах Unity төвтэй AR тоглоомын прототип системийг онолын үндэслэлтэйгээр зохиомжлон, хөгжүүлж, туршилтаар үнэлэхэд оршино.

Дээрх зорилгыг хэрэгжүүлэхийн тулд дараах үндсэн зорилтуудыг дэвшүүлэв.

\begin{enumerate}
    \item Өргөтгөсөн бодит байдал, байршилд суурилсан үйлчилгээ, тоглоомжуулалт, хэрэглэгчийн оролцооны онолын үндсийг судлан нэгтгэх.
    \item Мобайл AR аялал жуулчлалын системд тохирох хэрэглэгчийн урсгал, шаардлага, архитектур, өгөгдлийн загварыг боловсруулах.
    \item Unity клиент, Node.js/Express backend, PostgreSQL/Prisma өгөгдлийн сан дээр суурилсан прототип систем хэрэгжүүлэх.
    \item Радиуст суурилсан идэвхжилт, динамик AR байршуулалт, reward claim хамгаалалт, collection retrieval зэрэг үндсэн логикийг хэрэгжүүлэх.
    \item Прототипийг туршилтын сценариудаар үнэлж, инженерийн давуу тал, хязгаарлалт, цаашдын хөгжүүлэлтийн чиглэлийг тодорхойлох.
\end{enumerate}

\section{Судалгааны асуулт, шинэлэг тал ба практик ач холбогдол}

Энэхүү ажлыг дараах үндсэн судалгааны асуултууд чиглүүлж байна.

\begin{itemize}
    \item Байршлын триггер ба AR дүрслэлийг ямар архитектураар салган загварчилбал мобайл орчинд тогтвортой туршлага бий болох вэ?
    \item Тоглоомжуулалтын ямар механик Монголын аялал жуулчлалын соёлын контенттай хамгийн зохистой нийцэх вэ?
    \item Prototype түвшний системд ямар бизнес дүрэм, өгөгдлийн хязгаарлалт хэрэглэгчийн туршлагын найдвартай байдлыг тодорхойлох вэ?
\end{itemize}

Ажлын шинэлэг тал нь Монголын аялал жуулчлалын хэрэглээнд зориулсан AR туршлагыг зөвхөн визуал эффект бус, байршил, соёлын өгүүлэмж, reward collection, өгөгдлийн бүрэн бүтэн байдал, техникийн хязгаарлалтын огтлолцол дээр авч үзсэнд оршино. Өөрөөр хэлбэл энэ ажил нь ``AR үзүүлдэг аппликейшн'' хийхээс илүүтэйгээр байршлын логик ба орон зайн дүрслэлийг үүргээр нь салгаж, макро түвшинд GPS триггер, микро түвшинд AR plane placement ашиглах инженерийн шийдлийг академик үндэслэлтэй тайлбарлаж байгаа юм.

Практик ач холбогдлын хувьд уг прототип нь түүхэн дурсгал, музей, хотын маршрут, тусгай аяллын хөтөлбөр, соёлын арга хэмжээ зэрэгт дасан зохицох боломжтой суурь архитектур санал болгож байна. Систем нь цаашид олон байршил, улирлын контент, тусгай арга хэмжээ, analytics, админ удирдлага зэрэг өргөтгөлд шилжих боломжтой байдлаар загварчлагдсан.

\section{Судалгааны арга зүй}

Энэхүү дипломын ажлыг гүйцэтгэхдээ чанарын болон инженерийн холимог арга зүй ашиглав. Нэгдүгээрт, AR, LBS, тоглоомжуулалт, хэрэглэгчийн оролцооны чиглэлийн онолын болон хэрэглээний судалгааны материалуудыг харьцуулан шинжилж, системийн онолын хүрээг тогтоосон. Хоёрдугаарт, шаардлагын шинжилгээ болон хэрэглэгчийн урсгалын задлан шинжилгээнд тулгуурлан прототип системийн функциональ ба функциональ бус шаардлагыг тодорхойлов. Гуравдугаарт, Unity болон backend технологи ашиглан prototype хэрэгжүүлж, архитектурын шийдлүүдийг бодит төхөөрөмжийн нөхцөлтэй тулган үнэлсэн.

Туршилтын шатанд сценари-д суурилсан үнэлгээ ашигласан. Үүнд радиусын хил дээрх шилжилт, GPS савлалт, plane detection амжилтгүй тохиолдол, давхар reward claim оролдлого, сүлжээ тасалдсан үеийн fallback зэрэг бодит эрсдэл үүсгэх нөхцөлийг сонгож ажиглалт хийсэн. Ийнхүү энэхүү дипломын ажил нь зөвхөн тайлбарын түвшинд үлдээгүй, харин онол-зохиомж-хэрэгжилт-туршилтын гинжин холбоог хадгалсан инженерийн судалгааны хэлбэрээр хийгдсэн болно.

\section{Тайлангийн бүтэц}

Энэхүү дипломын тайлан нь үндсэндээ таван хэсэгтэй. Удиртгал бүлэгт асуудлын үндэслэл, зорилго, судалгааны асуулт, арга зүй, ажлын ач холбогдлыг тодорхойлов. Нэгдүгээр бүлэгт AR, LBS, тоглоомжуулалт, AR UX, архитектурын загвар, case study зэрэг онолын үндсийг авч үзэв. Хоёрдугаар бүлэгт системийн шаардлага, тоглоомын механик, архитектурын шийдэл, өгөгдлийн загвар, бизнес дүрмийг системийн зохиомжийн түвшинд задлан шинжлэв.

Гуравдугаар бүлэгт Unity төвтэй клиент, AR модуль, backend/API, өгөгдлийн сан, fallback шийдэл, туршилтын сценариудыг хэрэгжүүлэлтийн түвшинд тайлбарлав. Дүгнэлтэд ажлын үндсэн үр дүн, оруулсан хувь нэмэр, хязгаарлалт, цаашдын хөгжүүлэлтийн чиглэлийг нэгтгэн дүгнэв. Хавсралтад системийн гол кодын листинг болон өгөгдлийн сангийн схемийг оруулсан.

\end{sloppy}
